\documentclass[11pt,a4paper]{article}
\usepackage[utf8]{inputenc}
\usepackage[T1]{fontenc}
\usepackage{lmodern}
\usepackage{amsmath,amssymb,amsthm,mathtools}
\usepackage[margin=0.82in,headheight=14pt,top=0.85in,bottom=0.85in]{geometry}
\usepackage{enumitem}
\usepackage{fancyhdr}
\usepackage{xcolor}
\usepackage{mdframed}
\usepackage{hyperref}

% -------------------------------------------------------------------------
% Color Palette: Academic Executive Navy & Slate
% -------------------------------------------------------------------------
\definecolor{navyDark}{RGB}{18, 43, 76}
\definecolor{navyMid}{RGB}{28, 70, 120}
\definecolor{tealDark}{RGB}{14, 88, 82}
\definecolor{slateDark}{RGB}{55, 65, 81}
\definecolor{purpleDark}{RGB}{88, 44, 132}
\definecolor{amberDark}{RGB}{150, 95, 10}

\definecolor{bgThm}{RGB}{246, 249, 254}
\definecolor{bgDef}{RGB}{244, 250, 249}
\definecolor{bgEx}{RGB}{249, 250, 251}
\definecolor{bgProb}{RGB}{248, 246, 253}
\definecolor{bgRem}{RGB}{254, 253, 247}

% -------------------------------------------------------------------------
% mdframed Box Styles for Theorem Environments
% -------------------------------------------------------------------------
\newmdenv[
  backgroundcolor=bgThm,
  linecolor=navyMid,
  linewidth=2.5pt,
  topline=false,
  bottomline=false,
  rightline=false,
  innertopmargin=6pt,
  innerbottommargin=6pt,
  innerleftmargin=11pt,
  innerrightmargin=11pt,
  skipabove=6pt,
  skipbelow=6pt,
  nobreak=true
]{thmbox}

\newmdenv[
  backgroundcolor=bgDef,
  linecolor=tealDark,
  linewidth=2.5pt,
  topline=false,
  bottomline=false,
  rightline=false,
  innertopmargin=6pt,
  innerbottommargin=6pt,
  innerleftmargin=11pt,
  innerrightmargin=11pt,
  skipabove=6pt,
  skipbelow=6pt,
  nobreak=true
]{defbox}

\newmdenv[
  backgroundcolor=bgEx,
  linecolor=slateDark,
  linewidth=2.5pt,
  topline=false,
  bottomline=false,
  rightline=false,
  innertopmargin=6pt,
  innerbottommargin=6pt,
  innerleftmargin=11pt,
  innerrightmargin=11pt,
  skipabove=6pt,
  skipbelow=6pt,
  nobreak=true
]{exbox}

\newmdenv[
  backgroundcolor=bgProb,
  linecolor=purpleDark,
  linewidth=2.5pt,
  topline=false,
  bottomline=false,
  rightline=false,
  innertopmargin=6pt,
  innerbottommargin=6pt,
  innerleftmargin=11pt,
  innerrightmargin=11pt,
  skipabove=6pt,
  skipbelow=6pt,
  nobreak=true
]{probbox}

\newmdenv[
  backgroundcolor=bgRem,
  linecolor=amberDark,
  linewidth=2.5pt,
  topline=false,
  bottomline=false,
  rightline=false,
  innertopmargin=5pt,
  innerbottommargin=5pt,
  innerleftmargin=11pt,
  innerrightmargin=11pt,
  skipabove=6pt,
  skipbelow=6pt,
  nobreak=true
]{rembox}

% -------------------------------------------------------------------------
% Hyperref Setup
% -------------------------------------------------------------------------
\hypersetup{
    colorlinks=true,
    linkcolor=navyMid,
    citecolor=navyMid,
    urlcolor=tealDark,
    pdftitle={Lecture Notes: Group Actions and Ring Theory},
    pdfauthor={Abstract Algebra Lecture Series}
}

% -------------------------------------------------------------------------
% Header & Footer
% -------------------------------------------------------------------------
\pagestyle{fancy}
\fancyhf{}
\fancyhead[L]{\small\textsf{\textbf{\color{navyDark}Abstract Algebra: Group Actions \& Ring Theory}}}
\fancyhead[R]{\small\textsf{\color{slateDark}Page \thepage}}
\renewcommand{\headrulewidth}{0.5pt}
\renewcommand{\footrulewidth}{0pt}

% -------------------------------------------------------------------------
% Section Styling Commands
% -------------------------------------------------------------------------
\newcommand{\sectiontitle}[1]{%
  \vspace{1.4em}%
  \noindent{\Large\textbf{\color{navyDark}#1}}%
  \vspace{0.3em}\par\nopagebreak%
  {\color{navyMid}\hrule height 1pt}%
  \vspace{0.6em}\par\nopagebreak%
}

\newcommand{\subsectiontitle}[1]{%
  \vspace{1.0em}%
  \noindent{\large\textbf{\color{slateDark}#1}}%
  \vspace{0.25em}\par\nopagebreak%
}

% -------------------------------------------------------------------------
% Standardized Math Operators
% -------------------------------------------------------------------------
\DeclareMathOperator{\Sym}{Sym}
\DeclareMathOperator{\Cl}{Cl}
\DeclareMathOperator{\charac}{char}

% Key Formula Box
\newcommand{\keyformula}[1]{%
  \begin{center}%
    \fcolorbox{navyMid}{bgThm}{\parbox{0.92\linewidth}{\vspace{0.3em}\centering #1\vspace{0.3em}}}%
  \end{center}%
}

\begin{document}

% =========================================================================
% Title Header
% =========================================================================
\begin{center}
    {\huge\textbf{\color{navyDark}Abstract Algebra: Lecture Notes}}\\[0.5em]
    {\Large\textsf{\color{navyMid}Group Actions, the Class Equation, and Fundamentals of Ring Theory}}\\[0.8em]
    {\normalsize\textsf{\color{slateDark}Comprehensive Reference Covering Derived Subgroups, Actions, Conjugacy, Rings, Domains, and Fields}}
\end{center}
\vspace{0.5em}
{\color{navyDark}\hrule height 1.5pt}
\vspace{1.5em}

% =========================================================================
\sectiontitle{1. Preliminaries: Derived Subgroups of $S_4$ and $A_4$}
% =========================================================================

\begin{thmbox}
\textbf{Proposition 1.1.} \textit{The derived subgroup (commutator subgroup) of the symmetric group $S_4$ is the alternating group $A_4$, that is, $S_4' = A_4$.}
\end{thmbox}

\begin{proof}
Recall that for any two permutations $\sigma, \tau \in S_4$, the commutator is defined by $[\sigma, \tau] = \sigma \tau \sigma^{-1} \tau^{-1}$. Every cycle of length 3 can be expressed as a commutator of transpositions. Specifically, for distinct elements $a, b, c \in \{1, 2, 3, 4\}$, consider the transpositions $(a\,b)$ and $(b\,c)$:
\[
[(a\,b), (b\,c)] = (a\,b)(b\,c)(a\,b)^{-1}(b\,c)^{-1} = (a\,b)(b\,c)(a\,b)(b\,c) = (a\,c\,b) = (a\,b\,c)^2.
\]
Because $(a\,b\,c) = (a\,c\,b)^2$, every 3-cycle is a commutator in $S_4$. Since the set of all 3-cycles generates the alternating group $A_4$, it follows that:
\[
A_4 \subseteq S_4'.
\]
Conversely, since the quotient group $S_4 / A_4 \cong \mathbb{Z}_2$ is abelian, the universal property of commutator subgroups implies that $S_4' \subseteq A_4$. Combining both inclusions gives:
\[
S_4' = A_4. \qedhere
\]
\end{proof}

\begin{thmbox}
\textbf{Proposition 1.2.} \textit{The derived subgroup of the alternating group $A_4$ is the Klein four-group $V_4$, that is, $A_4' = V_4$.}
\end{thmbox}

\begin{proof}
Recall that the normal subgroups of $A_4$ are precisely:
\[
\{I\}, \quad V_4, \quad \text{and} \quad A_4,
\]
where $V_4 = \{I, (1\,2)(3\,4), (1\,3)(2\,4), (1\,4)(2\,3)\}$. 

Since $A_4$ is non-abelian, its derived subgroup must be non-trivial:
\[
A_4' \ne \{I\}.
\]
Now consider the quotient group $A_4 / V_4$. The order of this quotient is:
\[
|A_4 / V_4| = \frac{|A_4|}{|V_4|} = \frac{12}{4} = 3.
\]
Since any group of prime order is cyclic, $A_4 / V_4$ is cyclic and hence abelian. By the characterization of derived subgroups, if the quotient $A_4 / N$ is abelian, then $A_4' \subseteq N$. Setting $N = V_4$, we obtain:
\[
A_4' \subseteq V_4.
\]
Because $A_4'$ is a normal subgroup of $A_4$ contained in $V_4$, and $A_4' \ne \{I\}$, the only possibility among the normal subgroups of $A_4$ is:
\[
A_4' = V_4. \qedhere
\]
\end{proof}

% =========================================================================
\sectiontitle{2. Group Actions and Permutation Representations}
% =========================================================================

\begin{defbox}
\textbf{Definition 2.1 (Left Group Action).} Let $G$ be a group and let $A$ be a non-empty set. A \textbf{left group action} of $G$ on $A$ is a function
\[
* : G \times A \longrightarrow A, \quad (g, a) \longmapsto g * a
\]
satisfying the following two axioms:
\begin{enumerate}[label=\textbf{\arabic*)}]
    \item \textbf{Associativity (Compatibility):} $g_1 * (g_2 * a) = (g_1 g_2) * a$ for all $g_1, g_2 \in G$ and all $a \in A$.
    \item \textbf{Identity:} $e * a = a$ for all $a \in A$, where $e$ is the identity element of $G$.
\end{enumerate}
When such an action is given, $A$ is called a \textbf{$G$-set}, and we say that $G$ acts on $A$ from the left.
\end{defbox}

\begin{defbox}
\textbf{Definition 2.2 (Right Group Action).} A \textbf{right group action} of $G$ on $A$ is a mapping $* : A \times G \to A$ satisfying:
\begin{enumerate}[label=\textbf{\arabic*)}]
    \item $(a * g_1) * g_2 = a * (g_1 g_2)$ for all $g_1, g_2 \in G$ and all $a \in A$.
    \item $a * e = a$ for all $a \in A$.
\end{enumerate}
\end{defbox}

\subsectiontitle{Standard Examples of Group Actions}

\begin{exbox}
\textbf{Example 2.3 (Trivial Action).} Let $G$ be any group and $A \ne \emptyset$. Define $* : G \times A \to A$ by:
\[
g * a = a \quad \forall\, g \in G, \; \forall\, a \in A.
\]
Then for any $g_1, g_2 \in G$ and $a \in A$, we have $g_1 * (g_2 * a) = g_1 * a = a$ and $(g_1 g_2) * a = a$, so associativity holds. Moreover, $e * a = a$. Hence, $*$ is a valid group action, termed the \textbf{trivial action}.
\end{exbox}

\begin{exbox}
\textbf{Example 2.4 (Action of a Group on Itself by Left Translation).} Let $A = G$ and define $* : G \times A \to A$ by left group multiplication:
\[
g * a = ga \quad \forall\, g \in G, \; \forall\, a \in A (= G).
\]
Verifying the axioms:
\begin{align*}
    g_1 * (g_2 * a) &= g_1 * (g_2 a) = g_1(g_2 a) = (g_1 g_2)a = (g_1 g_2) * a, \\
    e * a &= ea = a.
\end{align*}
Therefore, $*$ is a group action, known as the \textbf{left regular action} or \textbf{left translation}.
\end{exbox}

\begin{exbox}
\textbf{Example 2.5 (Right Multiplication and Inversion).} If one attempts to define $* : G \times A \to A$ on $A = G$ by $g * a = ag$, then:
\[
g_1 * (g_2 * a) = g_1 * (ag_2) = (ag_2)g_1 = a(g_2 g_1).
\]
However, $(g_1 g_2) * a = a(g_1 g_2)$. Unless $G$ is abelian, $a(g_2 g_1) \ne a(g_1 g_2)$, so this does \textbf{not} define a left action.

To define a left action using right multiplication, define $* : G \times A \to A$ by:
\[
g * a = ag^{-1} \quad \forall\, g \in G, \; \forall\, a \in A (= G).
\]
Verifying the axioms:
\begin{align*}
    g_1 * (g_2 * a) &= g_1 * (ag_2^{-1}) = (ag_2^{-1})g_1^{-1} = a(g_2^{-1} g_1^{-1}) = a(g_1 g_2)^{-1} = (g_1 g_2) * a, \\
    e * a &= ae^{-1} = ae = a.
\end{align*}
Thus, this is a valid left action, called the \textbf{right regular action} (via inversion).
\end{exbox}

\begin{exbox}
\textbf{Example 2.6 (Action by Conjugation).} Let $A = G$ and define $* : G \times A \to A$ by conjugation:
\[
g * a = gag^{-1} \quad \forall\, g \in G, \; \forall\, a \in A (= G).
\]
Verifying the axioms:
\begin{align*}
    g_1 * (g_2 * a) &= g_1 * (g_2 a g_2^{-1}) = g_1(g_2 a g_2^{-1})g_1^{-1} = (g_1 g_2)a(g_1 g_2)^{-1} = (g_1 g_2) * a, \\
    e * a &= eae^{-1} = a.
\end{align*}
Hence, $*$ is a group action, known as the \textbf{action by conjugation}.
\end{exbox}

\begin{exbox}
\textbf{Example 2.7 (Action on Left Cosets).} Let $H \le G$ and let $A = G/H = \{xH \mid x \in G\}$ be the collection of all left cosets of $H$ in $G$. Define $* : G \times A \to A$ by:
\[
g * (xH) = (gx)H \quad \forall\, g \in G, \; \forall\, xH \in A.
\]
Verifying the axioms:
\begin{align*}
    g_1 * (g_2 * xH) &= g_1 * ((g_2 x)H) = (g_1(g_2 x))H = ((g_1 g_2)x)H = (g_1 g_2) * xH, \\
    e * (xH) &= (ex)H = xH.
\end{align*}
Thus, left multiplication on cosets constitutes a valid group action.
\end{exbox}

\begin{exbox}
\textbf{Example 2.8 (Conjugation on Subgroups).} Let $A = \{H \mid H \le G\}$ be the set of all subgroups of $G$. Define $* : G \times A \to A$ by:
\[
g * H = gHg^{-1} \quad \forall\, g \in G, \; \forall\, H \in A.
\]
Since $gHg^{-1}$ is always a subgroup of $G$, $*$ maps into $A$. Verifying the axioms:
\begin{align*}
    g_1 * (g_2 * H) &= g_1 * (g_2 H g_2^{-1}) = g_1(g_2 H g_2^{-1})g_1^{-1} = (g_1 g_2)H(g_1 g_2)^{-1} = (g_1 g_2) * H, \\
    e * H &= eHe^{-1} = H.
\end{align*}
Thus, $*$ is a valid group action of $G$ on the set of its subgroups.
\end{exbox}

\begin{rembox}
\textbf{Remark 2.9 (Well-Definedness Property).} If $* : G \times A \to A$ is a group action, then because $*$ is a function, whenever $x, y \in A$ satisfy $x = y$, it follows that $(g, x) = (g, y)$, and therefore:
\[
g * x = g * y \quad \forall\, g \in G.
\]
\end{rembox}

\subsectiontitle{Equivalence with Permutation Representations}

\begin{thmbox}
\textbf{Theorem 2.10 (Action--Homomorphism Correspondence).} \textit{Let $G$ be a group and $A$ a non-empty set. Every homomorphism $f : G \to \Sym(A)$ defines a group action of $G$ on $A$. Conversely, every group action of $G$ on $A$ induces a unique homomorphism $f : G \to \Sym(A)$.}
\end{thmbox}

\begin{proof}
$(\implies)$ Let $f : G \to \Sym(A)$ be a homomorphism. For each $g \in G$, let $\sigma_g = f(g)$, where $\sigma_g : A \to A$ is a bijection. Since $f$ is a homomorphism, for all $g_1, g_2 \in G$:
\[
f(g_1 g_2) = f(g_1) \circ f(g_2) \implies \sigma_{g_1 g_2} = \sigma_{g_1} \circ \sigma_{g_2}.
\]
Define $* : G \times A \to A$ by $g * a = \sigma_g(a)$ for all $a \in A$. Then:
\begin{align*}
    g_1 * (g_2 * a) &= g_1 * (\sigma_{g_2}(a)) = \sigma_{g_1}(\sigma_{g_2}(a)) = (\sigma_{g_1} \circ \sigma_{g_2})(a) = \sigma_{g_1 g_2}(a) = (g_1 g_2) * a.
\end{align*}
Furthermore, since $f$ is a homomorphism, $f(e) = I_A$ (the identity map on $A$), so:
\[
e * a = \sigma_e(a) = I_A(a) = a.
\]
Hence, $*$ is a group action of $G$ on $A$.

\medskip
\noindent $(\impliedby)$ Conversely, suppose $* : G \times A \to A$ is a group action. For each $g \in G$, define $\sigma_g : A \to A$ by:
\[
\sigma_g(a) = g * a \quad \forall\, a \in A.
\]
We first prove that $\sigma_g$ is a bijection, so that $\sigma_g \in \Sym(A)$:
\begin{itemize}
    \item \textbf{Injectivity:} Let $a_1, a_2 \in A$ such that $\sigma_g(a_1) = \sigma_g(a_2)$. Then $g * a_1 = g * a_2$. Acting on the left by $g^{-1}$:
    \[
    g^{-1} * (g * a_1) = g^{-1} * (g * a_2) \implies (g^{-1} g) * a_1 = (g^{-1} g) * a_2 \implies e * a_1 = e * a_2 \implies a_1 = a_2.
    \]
    Thus, $\sigma_g$ is one-to-one.
    \item \textbf{Surjectivity:} For any $a \in A$, observe that $g^{-1} * a \in A$, and:
    \[
    \sigma_g(g^{-1} * a) = g * (g^{-1} * a) = (g g^{-1}) * a = e * a = a.
    \]
    Thus, $\sigma_g$ is onto.
\end{itemize}
Therefore, $\sigma_g \in \Sym(A)$. Now define $f : G \to \Sym(A)$ by $f(g) = \sigma_g$. For any $g_1, g_2 \in G$ and $a \in A$:
\[
\sigma_{g_1 g_2}(a) = (g_1 g_2) * a = g_1 * (g_2 * a) = \sigma_{g_1}(\sigma_{g_2}(a)) = (\sigma_{g_1} \circ \sigma_{g_2})(a).
\]
Consequently, $\sigma_{g_1 g_2} = \sigma_{g_1} \circ \sigma_{g_2}$, which means $f(g_1 g_2) = f(g_1) \circ f(g_2)$. Hence $f$ is a homomorphism.
\end{proof}

\begin{defbox}
\textbf{Definition 2.11 (Permutation Representation).} The homomorphism $f : G \to \Sym(A)$ induced by a group action is called the \textbf{permutation representation} associated with the action.
\end{defbox}

\begin{thmbox}
\textbf{Corollary 2.12 (Cayley's Theorem).} \textit{Every group $G$ is isomorphic to a subgroup of its symmetric group $\Sym(G)$.}
\end{thmbox}

\begin{proof}
Consider the left regular action of $G$ on itself ($A = G$) defined by $g * a = ga$. By Theorem 2.10, this action induces a permutation representation:
\[
f : G \longrightarrow \Sym(G), \quad f(g) = \sigma_g, \quad \text{where } \sigma_g(a) = ga.
\]
We examine the kernel of $f$. Suppose $f(g_1) = f(g_2)$. Then $\sigma_{g_1} = \sigma_{g_2}$, which implies:
\[
\sigma_{g_1}(a) = \sigma_{g_2}(a) \quad \forall\, a \in G.
\]
Setting $a = e \in G$:
\[
g_1 e = g_2 e \implies g_1 = g_2.
\]
Thus, $f$ is injective, which implies $\ker f = \{e\}$. By the First Isomorphism Theorem for groups:
\[
G / \ker f \cong f(G) \implies G / \{e\} \cong f(G) \implies G \cong f(G).
\]
Since $f(G) \le \Sym(G)$, $G$ is isomorphic to the subgroup $f(G)$ of $\Sym(G)$.
\end{proof}

% =========================================================================
\sectiontitle{3. Kernel of an Action, Stabilizers, and Orbits}
% =========================================================================

\begin{defbox}
\textbf{Definition 3.1 (Kernel of an Action).} Let $* : G \times A \to A$ be a group action. The \textbf{kernel} of the action $*$ is defined as:
\[
\ker * = \{g \in G \mid g * a = a, \; \forall\, a \in A\}.
\]
That is, $\ker *$ consists of all elements of $G$ that fix every element of $A$.
\end{defbox}

\begin{thmbox}
\textbf{Theorem 3.2.} \textit{For any group action $* : G \times A \to A$, the kernel is a normal subgroup of $G$, $\ker * \trianglelefteq G$.}
\end{thmbox}

\begin{proof}
Since $e * a = a$ for all $a \in A$, $e \in \ker *$, so $\ker * \ne \emptyset$.
\begin{enumerate}[label=\textbf{\arabic*)}]
    \item \textbf{Closure under multiplication:} Let $g_1, g_2 \in \ker *$. Then for any $a \in A$:
    \[
    (g_1 g_2) * a = g_1 * (g_2 * a) = g_1 * a = a \implies g_1 g_2 \in \ker *.
    \]
    \item \textbf{Closure under inverses:} Let $g \in \ker *$, so $g * a = a$ for all $a \in A$. Then:
    \[
    g^{-1} * a = g^{-1} * (g * a) = (g^{-1} g) * a = e * a = a \implies g^{-1} \in \ker *.
    \]
\end{enumerate}
Thus, $\ker * \le G$. (Normality follows directly from Theorem 3.3 below since $\ker *$ is the kernel of a homomorphism.)
\end{proof}

\begin{thmbox}
\textbf{Theorem 3.3.} \textit{If $f : G \to \Sym(A)$ is the permutation representation corresponding to $*$, then:}
\keyformula{$\ker f = \ker *$}
\end{thmbox}

\begin{proof}
By definition of the kernel of a group homomorphism:
\begin{align*}
    \ker f &= \{g \in G \mid f(g) = I_A\} \\
    &= \{g \in G \mid \sigma_g = I_A\} \\
    &= \{g \in G \mid \sigma_g(a) = I_A(a), \; \forall\, a \in A\} \\
    &= \{g \in G \mid g * a = a, \; \forall\, a \in A\} \\
    &= \ker *. \qedhere
\end{align*}
\end{proof}

\subsectiontitle{Stabilizers and Orbits}

\begin{defbox}
\textbf{Definition 3.4 (Stabilizer).} Let $G$ act on a set $A$ and let $a \in A$. The \textbf{stabilizer} of $a$ in $G$ (also called the isotropy subgroup) is:
\[
G_a = \{g \in G \mid g * a = a\}.
\]
\end{defbox}

\begin{thmbox}
\textbf{Theorem 3.5.} \textit{For each $a \in A$, the stabilizer $G_a$ is a subgroup of $G$, that is, $G_a \le G$.}
\end{thmbox}

\begin{proof}
Since $e * a = a$, we have $e \in G_a$, so $G_a \ne \emptyset$.
\begin{enumerate}[label=\textbf{\arabic*)}]
    \item Let $g_1, g_2 \in G_a$. Then $g_1 * a = a$ and $g_2 * a = a$. Consequently:
    \[
    (g_1 g_2) * a = g_1 * (g_2 * a) = g_1 * a = a \implies g_1 g_2 \in G_a.
    \]
    \item Let $g \in G_a$, so $g * a = a$. Acting by $g^{-1}$:
    \[
    g^{-1} * a = g^{-1} * (g * a) = (g^{-1} g) * a = e * a = a \implies g^{-1} \in G_a.
    \]
\end{enumerate}
By the two-step subgroup criterion, $G_a \le G$.
\end{proof}

\begin{exbox}
\textbf{Example 3.6 (Stabilizer under Conjugation is the Normalizer).} If $G$ acts on itself by conjugation ($g * a = gag^{-1}$ for $a \in G$), then for any $a \in G$:
\begin{align*}
    G_a &= \{g \in G \mid g * a = a\} = \{g \in G \mid gag^{-1} = a\} \\
    &= \{g \in G \mid ga = ag\} = N(a),
\end{align*}
where $N(a)$ is the \textbf{normalizer} of $a$ in $G$.
\end{exbox}

\begin{defbox}
\textbf{Definition 3.7 (Orbit).} Let $G$ act on $A$ and let $a \in A$. The \textbf{orbit} of $a$ under $G$ is the set of all elements in $A$ to which $a$ can be moved by elements of $G$:
\[
G \cdot a = \{x \in A \mid g * a = x \text{ for some } g \in G\} = \{g * a \mid g \in G\}.
\]
\end{defbox}

\begin{thmbox}
\textbf{Theorem 3.8 (Orbit-Stabilizer Theorem).} \textit{Let $G$ be a group acting on a set $A$, and let $a \in A$. There exists a bijection between the orbit $G \cdot a$ and the set $G/G_a$ of all left cosets of $G_a$ in $G$.}
\end{thmbox}

\begin{proof}
Let $G \cdot a = \{g * a \mid g \in G\}$ and $G/G_a = \{g G_a \mid g \in G\}$. Define the mapping:
\[
\phi : G \cdot a \longrightarrow G/G_a, \quad \phi(g * a) = g G_a.
\]
We establish that $\phi$ is well-defined, injective, and surjective:
\begin{itemize}
    \item \textbf{Well-defined:} Suppose $g_1 * a = g_2 * a$. Then:
    \begin{align*}
        g_2^{-1} * (g_1 * a) = g_2^{-1} * (g_2 * a) &\implies (g_2^{-1} g_1) * a = e * a = a \\
        &\implies g_2^{-1} g_1 \in G_a \\
        &\implies g_1 G_a = g_2 G_a \\
        &\implies \phi(g_1 * a) = \phi(g_2 * a).
    \end{align*}
    Thus, $\phi$ is independent of the choice of representative.
    
    \item \textbf{One-to-one (Injective):} Suppose $\phi(g_1 * a) = \phi(g_2 * a)$. Then:
    \begin{align*}
        g_1 G_a = g_2 G_a &\implies g_2^{-1} g_1 \in G_a \\
        &\implies (g_2^{-1} g_1) * a = a \\
        &\implies g_2 * ((g_2^{-1} g_1) * a) = g_2 * a \\
        &\implies (g_2 g_2^{-1} g_1) * a = g_2 * a \\
        &\implies g_1 * a = g_2 * a.
    \end{align*}
    Thus $\phi$ is injective.
    
    \item \textbf{Onto (Surjective):} Let $t \in G/G_a$. Then $t = g G_a$ for some $g \in G$. For this $g$, the element $g * a \in G \cdot a$ satisfies $\phi(g * a) = g G_a = t$. Hence $\phi$ is surjective.
\end{itemize}
Therefore, $\phi$ is a bijection between $G \cdot a$ and $G/G_a$.
\end{proof}

\begin{thmbox}
\textbf{Corollary 3.9.} \textit{If $G$ is a finite group, then for any $a \in A$:}
\keyformula{$|G \cdot a| = [G : G_a] = \dfrac{|G|}{|G_a|} \implies \boxed{|G| = |G \cdot a| \cdot |G_a|}$}
\end{thmbox}

% =========================================================================
\sectiontitle{4. Conjugacy Classes and the Class Equation}
% =========================================================================

\begin{thmbox}
\textbf{Theorem 4.1 (Orbit Partitioning).} \textit{Let $* : G \times A \to A$ be a group action. Define a relation $\sim$ on $A$ by:
\[
a \sim b \iff a = g * b \quad \text{for some } g \in G.
\]
Then $\sim$ is an equivalence relation on $A$, and the equivalence classes are the orbits $G \cdot a$.}
\end{thmbox}

\begin{proof}
We check the three defining properties of an equivalence relation:
\begin{enumerate}[label=\textbf{\arabic*)}]
    \item \textbf{Reflexive:} Since $e \in G$, $a = e * a$, so $a \sim a$ for all $a \in A$.
    \item \textbf{Symmetric:} Let $a \sim b$. Then $a = g * b$ for some $g \in G$. Acting by $g^{-1}$:
    \[
    g^{-1} * a = g^{-1} * (g * b) = (g^{-1} g) * b = e * b = b \implies b = g^{-1} * a \implies b \sim a.
    \]
    \item \textbf{Transitive:} Let $a \sim b$ and $b \sim c$. Then there exist $g_1, g_2 \in G$ such that $a = g_1 * b$ and $b = g_2 * c$. Substituting gives:
    \[
    a = g_1 * (g_2 * c) = (g_1 g_2) * c.
    \]
    Since $g_1 g_2 \in G$, this implies $a \sim c$.
\end{enumerate}
Thus, $\sim$ is an equivalence relation. The equivalence class containing $a \in A$ is:
\[
\Cl(a) = \{b \in A \mid a \sim b\} = \{g * a \mid g \in G\} = G \cdot a.
\]
Since the equivalence classes partition $A$, we obtain:
\[
A = \bigcup_{a \in A} \Cl(a) = \bigcup_{a \in A} G \cdot a. \qedhere
\]
\end{proof}

\begin{defbox}
\textbf{Definition 4.2 (Conjugacy Class).} Under the action of $G$ on itself by conjugation ($g * a = gag^{-1}$), the orbit of $a \in G$ is called the \textbf{conjugacy class} of $a$, denoted $\Cl(a)$:
\[
\Cl(a) = \{gag^{-1} \mid g \in G\}.
\]
\end{defbox}

\begin{probbox}
\textbf{Exercise 4.3 (Conjugacy as an Equivalence Relation).} \textit{Let $G$ be a group. Two elements $a, b \in G$ are said to be conjugate (written $a \sim b$) if there exists $c \in G$ such that $b = c a c^{-1}$. Prove directly that $\sim$ is an equivalence relation on $G$.}
\end{probbox}

\begin{proof}[Solution]
We verify the three equivalence relation axioms:
\begin{enumerate}[label=\textbf{\arabic*)}]
    \item \textbf{Reflexivity:} For any $a \in G$, $a = e a e^{-1}$ (since $e \in G$), so $a \sim a$.
    \item \textbf{Symmetry:} Let $a \sim b$. Then there exists $c \in G$ such that $a = c b c^{-1}$. Multiplying on the left by $c^{-1}$ and on the right by $c$:
    \[
    c^{-1} a c = c^{-1}(c b c^{-1})c = (c^{-1} c) b (c^{-1} c) = e b e = b.
    \]
    Writing $c = (c^{-1})^{-1}$, this gives $b = c^{-1} a (c^{-1})^{-1}$. Since $c^{-1} \in G$, we have $b \sim a$.
    \item \textbf{Transitivity:} Let $a \sim b$ and $b \sim c$. Then there exist $d_1, d_2 \in G$ such that:
    \[
    a = d_1 b d_1^{-1} \quad \text{and} \quad b = d_2 c d_2^{-1}.
    \]
    Substituting $b$ into the expression for $a$:
    \[
    a = d_1 (d_2 c d_2^{-1}) d_1^{-1} = (d_1 d_2) c (d_2^{-1} d_1^{-1}) = (d_1 d_2) c (d_1 d_2)^{-1}.
    \]
    Since $d_1 d_2 \in G$, it follows that $a \sim c$.
\end{enumerate}
Thus, conjugacy is an equivalence relation on $G$.
\end{proof}

\begin{thmbox}
\textbf{Proposition 4.4 (Conjugacy Classes and the Center).} \textit{Let $G$ be a group and let $a \in G$. Then:}
\begin{enumerate}[label=\textbf{\arabic*)}]
    \item $\Cl(a) = \{a\} \iff a \in Z(G)$.
    \item $N(a) = G \iff a \in Z(G)$.
\end{enumerate}
\end{thmbox}

\begin{proof}
\textbf{Part 1:} 
$(\implies)$ Let $\Cl(a) = \{a\}$. For any $g \in G$, $gag^{-1} \in \Cl(a) = \{a\}$, so $gag^{-1} = a$. Multiplying on the right by $g$ yields $ga = ag$. Since $g$ was arbitrary, $a$ commutes with every element of $G$, so $a \in Z(G)$.

\noindent $(\impliedby)$ Conversely, suppose $a \in Z(G)$. Then for all $g \in G$, $ga = ag$, so $gag^{-1} = a$. Thus:
\[
\Cl(a) = \{gag^{-1} \mid g \in G\} = \{agg^{-1} \mid g \in G\} = \{a\}.
\]

\medskip
\noindent \textbf{Part 2:}
$(\implies)$ Suppose $N(a) = G$. Then every $g \in G$ satisfies $g \in N(a)$, meaning $gag^{-1} = a$, or equivalently $ga = ag$. Thus $a \in Z(G)$.

\noindent $(\impliedby)$ Conversely, if $a \in Z(G)$, then $ga = ag$ for all $g \in G$, which means $g \in N(a)$ for all $g \in G$. Hence $N(a) = G$.
\end{proof}

\subsectiontitle{The Class Equation}

For any finite group $G$, because the distinct conjugacy classes partition $G$, we have:
\[
|G| = \sum_{a \in G} |\Cl(a)|.
\]
By the Orbit-Stabilizer Theorem, $|\Cl(a)| = [G : N(a)] = \frac{|G|}{|N(a)|}$. Separating the singleton conjugacy classes (which correspond precisely to elements of the center $Z(G)$) from the non-central classes yields:

\keyformula{
\textbf{The Class Equation}\\[0.4em]
$\displaystyle |G| = |Z(G)| + \sum_{a \notin Z(G)} \left|\frac{G}{N(a)}\right| = |Z(G)| + \sum_{a \notin Z(G)} [G : N(a)]$
}
where the sum runs over one representative $a$ from each conjugacy class not contained in $Z(G)$.

\subsectiontitle{Applications of the Class Equation to $p$-Groups}

\begin{thmbox}
\textbf{Theorem 4.5 (Center of a $p$-Group is Non-Trivial).} \textit{Let $G$ be a finite group of order $|G| = p^n$, where $p$ is a prime and $n \ge 1$. Then $|Z(G)| \ge p$, that is, the center of a non-trivial $p$-group is non-trivial.}
\end{thmbox}

\begin{proof}
Let $|Z(G)| = r$. By the class equation:
\[
|G| = |Z(G)| + \sum_{a \notin Z(G)} [G : N(a)]. \tag*{(A)}
\]
For any $a \notin Z(G)$, the normalizer $N(a)$ is a proper subgroup of $G$, so $N(a) \lneq G$. By Lagrange's Theorem, $|N(a)|$ divides $|G| = p^n$, so $|N(a)| = p^{\alpha_a}$ with $\alpha_a < n$. Therefore:
\[
[G : N(a)] = \frac{|G|}{|N(a)|} = \frac{p^n}{p^{\alpha_a}} = p^{n - \alpha_a}, \quad \text{where } n - \alpha_a \ge 1.
\]
Thus, $p \mid [G : N(a)]$ for each term in the summation, which implies:
\[
p \;\Big|\; \sum_{a \notin Z(G)} [G : N(a)].
\]
Since $|G| = p^n$, we also have $p \mid |G|$. Rewriting the class equation as:
\[
|Z(G)| = |G| - \sum_{a \notin Z(G)} [G : N(a)],
\]
it follows that $p \mid |Z(G)|$. Since the identity element $e \in Z(G)$, the center is non-empty, so $|Z(G)| \ge 1$. Because $|Z(G)|$ is a positive integer divisible by the prime $p$, we conclude that:
\[
|Z(G)| \ge p. \qedhere
\]
\end{proof}

\begin{thmbox}
\textbf{Theorem 4.6 (Groups of Order $p^2$).} \textit{Every group $G$ of order $p^2$, where $p$ is a prime, is abelian.}
\end{thmbox}

\begin{proof}
By Theorem 4.5, $|Z(G)| \ge p$, and by Lagrange's Theorem, $|Z(G)|$ must divide $|G| = p^2$. Hence:
\[
|Z(G)| = p \quad \text{or} \quad |Z(G)| = p^2.
\]
Suppose for contradiction that $|Z(G)| = p$. Then the quotient group has order:
\[
|G / Z(G)| = \frac{|G|}{|Z(G)|} = \frac{p^2}{p} = p.
\]
Since any group of prime order is cyclic, $G / Z(G)$ is cyclic. It is a standard result that if $G / Z(G)$ is cyclic, then $G$ is abelian. But if $G$ is abelian, then $Z(G) = G$, so $|Z(G)| = p^2$, which contradicts $|Z(G)| = p$.

Hence, $|Z(G)| \ne p$. The only remaining possibility is $|Z(G)| = p^2 = |G|$, which implies $Z(G) = G$. Therefore, $G$ is abelian.
\end{proof}

\begin{rembox}
\textbf{Remark 4.7.} If $H$ is a subgroup of $G$ such that $H \le Z(G)$, then $H$ is normal in $G$, $H \trianglelefteq G$. This holds because every element of $H$ commutes with all elements of $G$, so $g h g^{-1} = h \in H$ for all $g \in G, h \in H$.
\end{rembox}

\subsectiontitle{Conjugacy Classes in Symmetric Groups $S_n$}

Two permutations in $S_n$ are conjugate if and only if they have the same cycle decomposition structure (cycle type).

\begin{thmbox}
\textbf{Theorem 4.8 (Size of a Conjugacy Class in $S_n$).} \textit{Let $\sigma \in S_n$ have cycle type consisting of $k_i$ cycles of length $m_i$ for $i = 1, 2, \dots, r$ (including 1-cycles), such that $\sum_{i=1}^r k_i m_i = n$. Then the number of conjugates of $\sigma$ in $S_n$ is:}
\keyformula{$\displaystyle |\Cl(\sigma)| = \frac{n!}{(k_1! \, m_1^{k_1})(k_2! \, m_2^{k_2}) \cdots (k_r! \, m_r^{k_r})}$}
\end{thmbox}

\begin{exbox}
\textbf{Example 4.9 (Calculations in $S_4$).} Here $n = 4$, so $4! = 24$.
\begin{itemize}
    \item \textbf{Transpositions (cycles of length 2):} The cycle type is $2^1 1^2$ ($k_1 = 1, m_1 = 2$ and $k_2 = 2, m_2 = 1$).
    \[
    |\Cl((1\,2))| = \frac{4!}{(1! \cdot 2^1)(2! \cdot 1^2)} = \frac{24}{2 \times 2} = 6.
    \]
    Explicitly: $\Cl((1\,2)) = \{(1\,2), (1\,3), (1\,4), (2\,3), (2\,4), (3\,4)\}$.
    \item \textbf{3-cycles:} The cycle type is $3^1 1^1$ ($k_1 = 1, m_1 = 3$ and $k_2 = 1, m_2 = 1$).
    \[
    |\Cl((1\,2\,3))| = \frac{4!}{(1! \cdot 3^1)(1! \cdot 1^1)} = \frac{24}{3 \times 1} = 8.
    \]
\end{itemize}
\end{exbox}

\begin{exbox}
\textbf{Example 4.10 (Calculations in $S_5$).} Here $n = 5$, so $5! = 120$.
\begin{itemize}
    \item \textbf{Transpositions (cycles of length 2):} The cycle type is $2^1 1^3$ ($k_1 = 1, m_1 = 2$ and $k_2 = 3, m_2 = 1$).
    \[
    |\Cl((1\,2))| = \frac{5!}{(1! \cdot 2^1)(3! \cdot 1^3)} = \frac{120}{2 \times 6} = \frac{120}{12} = 10.
    \]
\end{itemize}
\end{exbox}

% =========================================================================
\sectiontitle{5. Introduction to Ring Theory}
% =========================================================================

\begin{defbox}
\textbf{Definition 5.1 (Ring).} A mathematical system $(R, +, \cdot)$ consisting of a non-empty set $R$ equipped with two binary operations `$+$' (addition) and `$\cdot$' (multiplication) is called a \textbf{ring} if it satisfies:
\begin{enumerate}[label=\textbf{\arabic*)}]
    \item \textbf{$(R, +)$ is an abelian group:}
    \begin{itemize}
        \item Closure: $a + b \in R$ for all $a, b \in R$.
        \item Associativity: $(a + b) + c = a + (b + c)$ for all $a, b, c \in R$.
        \item Additive identity: There exists an element $0 \in R$ such that $a + 0 = 0 + a = a$ for all $a \in R$.
        \item Additive inverse: For each $a \in R$, there exists $-a \in R$ such that $a + (-a) = (-a) + a = 0$.
        \item Commutativity: $a + b = b + a$ for all $a, b \in R$.
    \end{itemize}
    \item \textbf{$(R, \cdot)$ is a semigroup:}
    \begin{itemize}
        \item Closure: $a \cdot b \in R$ for all $a, b \in R$.
        \item Associativity: $(a \cdot b) \cdot c = a \cdot (b \cdot c)$ for all $a, b, c \in R$.
    \end{itemize}
    \item \textbf{Distributivity of multiplication over addition:} For all $a, b, c \in R$:
    \begin{align*}
        a \cdot (b + c) &= a \cdot b + a \cdot c \quad \text{(Left distributive law)}, \\
        (a + b) \cdot c &= a \cdot c + b \cdot c \quad \text{(Right distributive law)}.
    \end{align*}
\end{enumerate}
\end{defbox}

\begin{defbox}
\textbf{Definition 5.2 (Special Classes of Rings).}
\begin{enumerate}[label=\textbf{\arabic*)}]
    \item \textbf{Zero Element:} The identity of $(R, +)$ is called the \textbf{zero element} of the ring, denoted $0$.
    \item \textbf{Ring with Unity:} If $(R, \cdot)$ has an identity element $1 \in R$ such that $a \cdot 1 = 1 \cdot a = a$ for all $a \in R$, then $R$ is called a \textbf{ring with unity} (or identity).
    \item \textbf{Commutative Ring:} If multiplication is commutative ($a \cdot b = b \cdot a$ for all $a, b \in R$), then $R$ is called a \textbf{commutative ring}.
\end{enumerate}
\end{defbox}

\subsectiontitle{Twelve Standard Examples of Rings}

\begin{enumerate}[label=\textbf{\arabic*)}]
    \item \textbf{Number Systems:} The systems $(\mathbb{Z}, +, \cdot)$, $(\mathbb{Q}, +, \cdot)$, $(\mathbb{R}, +, \cdot)$, and $(\mathbb{C}, +, \cdot)$ under standard addition and multiplication are all commutative rings with unity $1$.
    
    \item \textbf{Integers Modulo $n$:} For any integer $n \ge 2$, $(\mathbb{Z}_n, +_n, \times_n)$ is a commutative ring with unity $[1]$.
    
    \item \textbf{Quadratic Ring of Integers:} For a square-free integer $p$, $\mathbb{Z}[\sqrt{p}] = \{a + b\sqrt{p} \mid a, b \in \mathbb{Z}\}$ forms a commutative ring with unity $1$.
    
    \item \textbf{Ring of Gaussian Integers:} The set $\mathbb{Z}[i] = \{a + bi \mid a, b \in \mathbb{Z}\}$ under complex addition and multiplication is a commutative ring with unity $1$.
    
    \item \textbf{Real Matrix Ring:} For $n \ge 2$, $M_n(\mathbb{R})$ (the set of all $n \times n$ real matrices) under matrix addition and matrix multiplication satisfies:
    \begin{itemize}
        \item $(M_n(\mathbb{R}), +)$ is an abelian group with zero matrix $0_{n \times n}$.
        \item $(M_n(\mathbb{R}), \cdot)$ is a semigroup with identity $I_n$.
        \item Matrix multiplication distributes over matrix addition: $A(B + C) = AB + AC$ and $(A + B)C = AC + BC$.
    \end{itemize}
    Since matrix multiplication is generally non-commutative, the ring $(M_n(\mathbb{R}), +, \cdot)$ is a \textbf{non-commutative ring with unity}.
    
    \item \textbf{Integer Matrix Ring:} $M_n(\mathbb{Z})$ is similarly a non-commutative ring with unity $I_n$.
    
    \item \textbf{A Constant $2 \times 2$ Matrix Ring with Non-Standard Unity:} Consider the set:
    \[
    A = \left\{ \begin{bmatrix} x & x \\ x & x \end{bmatrix} \;\middle|\; x \in \mathbb{R} \right\}.
    \]
    Under standard matrix addition, $(A, +)$ is an abelian group. Under matrix multiplication:
    \[
    \begin{bmatrix} x & x \\ x & x \end{bmatrix} \begin{bmatrix} y & y \\ y & y \end{bmatrix} = \begin{bmatrix} 2xy & 2xy \\ 2xy & 2xy \end{bmatrix} = \begin{bmatrix} y & y \\ y & y \end{bmatrix} \begin{bmatrix} x & x \\ x & x \end{bmatrix}.
    \]
    Thus, multiplication is commutative on $A$. To find the unity $E = \begin{bmatrix} e & e \\ e & e \end{bmatrix} \in A$:
    \[
    \begin{bmatrix} x & x \\ x & x \end{bmatrix} \begin{bmatrix} e & e \\ e & e \end{bmatrix} = \begin{bmatrix} x & x \\ x & x \end{bmatrix} \implies \begin{bmatrix} 2xe & 2xe \\ 2xe & 2xe \end{bmatrix} = \begin{bmatrix} x & x \\ x & x \end{bmatrix} \implies 2xe = x \implies e = \frac{1}{2}.
    \]
    Hence, the unity of the ring $A$ is:
    \[
    E = \begin{bmatrix} 1/2 & 1/2 \\ 1/2 & 1/2 \end{bmatrix}.
    \]
    Furthermore, for any non-zero element $X = \begin{bmatrix} x & x \\ x & x \end{bmatrix}$ ($x \ne 0$), its multiplicative inverse $Y = \begin{bmatrix} y & y \\ y & y \end{bmatrix}$ satisfies:
    \[
    XY = E \implies 2xy = \frac{1}{2} \implies y = \frac{1}{4x} \implies Y = \begin{bmatrix} \frac{1}{4x} & \frac{1}{4x} \\[0.3em] \frac{1}{4x} & \frac{1}{4x} \end{bmatrix}.
    \]
    Thus $(A \setminus \{0\}, \cdot)$ is an abelian group, making $A$ a field!
    
    \item \textbf{Matrix Ring with Projection Unity:} The set:
    \[
    A = \left\{ \begin{bmatrix} x & 0 \\ 0 & 0 \end{bmatrix} \;\middle|\; x \in \mathbb{R} \right\}
    \]
    is a commutative ring with unity $E = \begin{bmatrix} 1 & 0 \\ 0 & 0 \end{bmatrix}$.
    
    \item \textbf{The Trivial Ring:} The singleton set $R = \{0\}$ with $0 + 0 = 0$ and $0 \cdot 0 = 0$ is a ring with unity $0$, called the \textbf{trivial ring} (or \textbf{zero ring}). Any ring with $R \ne \{0\}$ is called a \textbf{non-trivial ring}.
    
    \item \textbf{Polynomial Ring over $\mathbb{Z}$:} The set $\mathbb{Z}[x]$ of all polynomials with integer coefficients is a commutative ring with unity 1.
    
    \item \textbf{Polynomial Rings over Fields:} The sets $\mathbb{Q}[x], \mathbb{R}[x]$, and $\mathbb{C}[x]$ are all commutative rings with unity 1.
    
    \item \textbf{Ring of Continuous Functions:} Let $C[a, b]$ denote the set of all continuous real-valued functions defined on the closed interval $[a, b]$. Under pointwise addition and multiplication:
    \[
    (f + g)(x) = f(x) + g(x), \qquad (f \cdot g)(x) = f(x) \cdot g(x) \quad \forall\, x \in [a, b],
    \]
    $C[a, b]$ is a commutative ring with unity $e(x) \equiv 1$ for all $x \in [a, b]$.
\end{enumerate}

\subsectiontitle{Fundamental Algebraic Properties of a Ring}

\begin{thmbox}
\textbf{Theorem 5.3 (Ring Identities).} \textit{Let $(R, +, \cdot)$ be a ring. Then for all $a, b, c \in R$:}
\begin{enumerate}[label=\textbf{\arabic*)}]
    \item $a \cdot 0 = 0 \cdot a = 0$
    \item $a(-b) = (-a)b = -(ab)$
    \item $(-a)(-b) = ab$
    \item $a(b - c) = ab - ac$
    \item $(b - c)a = ba - ca$
\end{enumerate}
\end{thmbox}

\begin{proof}
\textbf{Part 1:} In $(R, +)$, $0 + 0 = 0$. By left distributivity:
\[
a \cdot 0 = a \cdot (0 + 0) = a \cdot 0 + a \cdot 0.
\]
Adding the additive identity $0$ to the left-hand side gives $0 + a \cdot 0 = a \cdot 0 + a \cdot 0$. By the left cancellation law in the abelian group $(R, +)$:
\[
a \cdot 0 = 0.
\]
Similarly, $(0 + 0) \cdot a = 0 \cdot a + 0 \cdot a \implies 0 \cdot a = 0$.

\medskip
\noindent \textbf{Part 2:} By distributivity and Part 1:
\[
a(-b) + ab = a(-b + b) = a \cdot 0 = 0.
\]
Since the additive inverse of $ab$ in $(R, +)$ is unique and denoted $-(ab)$, it follows that $a(-b) = -(ab)$. Similarly:
\[
(-a)b + ab = (-a + a)b = 0 \cdot b = 0 \implies (-a)b = -(ab).
\]

\medskip
\noindent \textbf{Part 3:} Applying Part 2 twice:
\[
(-a)(-b) = -((-a)b) = -(-(ab)) = ab.
\]

\medskip
\noindent \textbf{Part 4:} Recalling that $b - c = b + (-c)$ and using distributivity and Part 2:
\[
a(b - c) = a(b + (-c)) = ab + a(-c) = ab + (-(ac)) = ab - ac.
\]

\medskip
\noindent \textbf{Part 5:} Similarly, by right distributivity:
\[
(b - c)a = (b + (-c))a = ba + (-c)a = ba + (-(ca)) = ba - ca. \qedhere
\]
\end{proof}

% =========================================================================
\sectiontitle{6. Zero Divisors and Integral Domains}
% =========================================================================

\begin{defbox}
\textbf{Definition 6.1 (Zero Divisors).} Let $R$ be a ring. A non-zero element $a \in R$ is called a \textbf{zero divisor} (or divisor of zero) if there exists a non-zero element $b \in R$ such that:
\[
ab = 0 \quad \text{or} \quad ba = 0.
\]
More specifically:
\begin{itemize}
    \item $a \ne 0$ is a \textbf{left zero divisor} if there exists $b \ne 0$ such that $ab = 0$.
    \item $b \ne 0$ is a \textbf{right zero divisor} if there exists $a \ne 0$ such that $ab = 0$.
\end{itemize}
\end{defbox}

\begin{exbox}
\textbf{Example 6.2 (Zero Divisors in $\mathbb{Z}_6$).} In the ring $(\mathbb{Z}_6, +_6, \times_6)$, the non-zero elements $[2], [3], [4]$ are zero divisors because:
\[
[2] \times_6 [3] = [0], \qquad [3] \times_6 [4] = [0].
\]
In contrast, in the rings $\mathbb{Z}, \mathbb{Q}, \mathbb{R}$, and $\mathbb{C}$, there are no zero divisors.
\end{exbox}

\begin{defbox}
\textbf{Definition 6.3 (Integral Domain).} A commutative ring $R$ with unity ($1 \ne 0$) is called an \textbf{integral domain} if it has no zero divisors, that is, for all $a, b \in R$:
\[
ab = 0 \implies a = 0 \quad \text{or} \quad b = 0.
\]
\end{defbox}

\begin{exbox}
\textbf{Examples and Non-Examples of Integral Domains.}
\begin{enumerate}[label=\textbf{\arabic*)}]
    \item \textbf{Examples:}
    \begin{itemize}
        \item The number systems $(\mathbb{Z}, +, \cdot)$, $(\mathbb{Q}, +, \cdot)$, $(\mathbb{R}, +, \cdot)$, and $(\mathbb{C}, +, \cdot)$ are integral domains.
        \item $(\mathbb{Z}_p, +_p, \times_p)$ is an integral domain if and only if $p$ is a prime.
        \item The polynomial rings $\mathbb{Z}[x], \mathbb{Q}[x], \mathbb{R}[x]$, and $\mathbb{C}[x]$ are integral domains.
        \item The Gaussian integers $\mathbb{Z}[i]$ form an integral domain.
    \end{itemize}
    \item \textbf{Non-Examples:}
    \begin{itemize}
        \item For composite $n$, $(\mathbb{Z}_n, +_n, \times_n)$ is not an integral domain (for example, in $\mathbb{Z}_4$, $[2] \times_4 [2] = [0]$ with $[2] \ne [0]$).
        \item For $n \ge 2$, $M_n(\mathbb{R})$ is not an integral domain because it is non-commutative and possesses zero divisors (e.g., $\begin{bmatrix} 0 & 1 \\ 0 & 0 \end{bmatrix}^2 = 0$).
    \end{itemize}
\end{enumerate}
\end{exbox}

\begin{thmbox}
\textbf{Theorem 6.4 (Cancellation Law in Integral Domains).} \textit{A commutative ring $R$ is an integral domain if and only if the cancellation law holds under multiplication, that is, for all $a, b, c \in R$ with $a \ne 0$:}
\[
ab = ac \implies b = c.
\]
\end{thmbox}

\begin{proof}
$(\implies)$ Let $R$ be an integral domain, and let $ab = ac$ with $a \ne 0$. Then:
\[
ab - ac = 0 \implies a(b - c) = 0.
\]
Since $R$ is an integral domain and $a \ne 0$, $R$ contains no zero divisors, which forces $b - c = 0$, so $b = c$.

\medskip
\noindent $(\impliedby)$ Conversely, assume that the cancellation law holds in $R$. Suppose $ab = 0$ with $a \ne 0$. By Theorem 5.3, $a \cdot 0 = 0$, so:
\[
ab = a \cdot 0.
\]
Since $a \ne 0$, the cancellation law immediately gives $b = 0$. Therefore, $R$ has no zero divisors, so $R$ is an integral domain.
\end{proof}

% =========================================================================
\sectiontitle{7. Fields, Division Rings, and Ring Characteristic}
% =========================================================================

\begin{defbox}
\textbf{Definition 7.1 (Field).} A commutative ring $(F, +, \cdot)$ with unity ($1 \ne 0$) is called a \textbf{field} if every non-zero element has a multiplicative inverse in $F$, that is:
\[
\forall\, a \in F \setminus \{0\}, \quad \exists\, a^{-1} \in F \quad \text{such that} \quad a \cdot a^{-1} = a^{-1} \cdot a = 1.
\]
Equivalently, $(F, +, \cdot)$ is a field if:
\begin{enumerate}[label=\textbf{\arabic*)}]
    \item $(F, +)$ is an abelian group.
    \item $(F \setminus \{0\}, \cdot)$ is an abelian group.
    \item Distributive laws hold.
\end{enumerate}
\end{defbox}

\begin{defbox}
\textbf{Definition 7.2 (Division Ring / Skew-Field).} A ring $(R, +, \cdot)$ with unity ($1 \ne 0$) is called a \textbf{division ring} (or \textbf{skew-field}) if every non-zero element has a multiplicative inverse. (Multiplication need not be commutative.)
\end{defbox}

\begin{exbox}
\textbf{Examples and Non-Examples of Fields.}
\begin{itemize}
    \item $(\mathbb{Q}, +, \cdot)$, $(\mathbb{R}, +, \cdot)$, and $(\mathbb{C}, +, \cdot)$ are fields.
    \item $(\mathbb{Z}_p, +_p, \times_p)$ is a field if and only if $p$ is prime.
    \item $(\mathbb{Z}, +, \cdot)$ is \textbf{not} a field because only $\pm 1$ have multiplicative inverses in $\mathbb{Z}$.
    \item $(\mathbb{Z}_n, +_n, \times_n)$ for composite $n$ is not a field.
\end{itemize}
\end{exbox}

\begin{thmbox}
\textbf{Theorem 7.3.} \textit{Every field is an integral domain.}
\end{thmbox}

\begin{proof}
Let $F$ be a field. Then $F$ is a commutative ring with unity $1 \ne 0$. To prove that $F$ is an integral domain, we must show that it has no zero divisors.

Let $a, b \in F$ such that $ab = 0$. If $a = 0$, we are done. If $a \ne 0$, then because $F$ is a field, the multiplicative inverse $a^{-1} \in F$ exists. Multiplying both sides by $a^{-1}$:
\begin{align*}
    a^{-1}(ab) &= a^{-1} \cdot 0 \\
    (a^{-1}a)b &= 0 \\
    1 \cdot b &= 0 \\
    b &= 0.
\end{align*}
Thus, $ab = 0$ implies $a = 0$ or $b = 0$. Therefore, $F$ has no zero divisors, so $F$ is an integral domain.
\end{proof}

\begin{rembox}
\textbf{Remark 7.4.} The converse of Theorem 7.3 does not hold in general: $(\mathbb{Z}, +, \cdot)$ is an integral domain, but it is not a field. However, in the finite setting, the converse is universally true!
\end{rembox}

\begin{thmbox}
\textbf{Theorem 7.5 (Finite Integral Domains are Fields).} \textit{Every finite integral domain is a field.}
\end{thmbox}

\begin{proof}
Let $D$ be a finite integral domain. Since $D$ is finite, its elements can be listed as:
\[
D = \{x_1, x_2, \dots, x_n\},
\]
where $x_1, \dots, x_n$ are distinct elements. Let $a \in D$ be any non-zero element ($a \ne 0$). Consider the set of products:
\[
aD = \{ax_1, ax_2, \dots, ax_n\}.
\]
Since $D$ is closed under multiplication, $aD \subseteq D$. We claim that all $n$ elements of $aD$ are distinct:
\[
ax_i = ax_j \implies ax_i - ax_j = 0 \implies a(x_i - x_j) = 0.
\]
Because $D$ is an integral domain and $a \ne 0$, $D$ contains no zero divisors, which forces $x_i - x_j = 0$, so $x_i = x_j$. 

Hence, if $i \ne j$, then $ax_i \ne ax_j$. Therefore, $aD$ contains $n$ distinct elements of $D$. Since $D$ has exactly $n$ elements, by the Pigeonhole Principle:
\[
aD = D.
\]
Since $D$ is an integral domain with unity $1$, we have $1 \in D = aD$. Thus, there exists some $x_k \in D$ such that:
\[
ax_k = 1.
\]
Because $D$ is commutative, $ax_k = x_k a = 1$, which means $x_k = a^{-1}$. Thus, every non-zero element of $D$ possesses a multiplicative inverse in $D$. Hence, $D$ is a field.
\end{proof}

\begin{thmbox}
\textbf{Corollary 7.6.} \textit{$(\mathbb{Z}_p, +_p, \times_p)$ is a field if and only if $p$ is a prime number.}
\end{thmbox}

\subsectiontitle{Characteristic of a Ring}

\begin{defbox}
\textbf{Definition 7.7 (Characteristic).} Let $(R, +, \cdot)$ be a ring. If there exists a least positive integer $n$ such that:
\[
n \cdot a = \underbrace{a + a + \dots + a}_{n \text{ times}} = 0 \quad \forall\, a \in R,
\]
then $n$ is called the \textbf{characteristic} of the ring $R$, denoted $\charac(R) = n$. If no such positive integer exists, $R$ is said to have \textbf{characteristic zero} (or infinite characteristic), denoted $\charac(R) = 0$.
\end{defbox}

\begin{exbox}
\textbf{Examples of Characteristics.}
\begin{itemize}
    \item $\charac(\mathbb{Z}) = 0, \quad \charac(\mathbb{Q}) = 0, \quad \charac(\mathbb{R}) = 0, \quad \charac(\mathbb{C}) = 0$.
    \item $\charac(\mathbb{Z}_n) = n$ for any $n \ge 2$.
    \item $\charac(\mathbb{Z}_p) = p$ for any prime $p$.
\end{itemize}
\end{exbox}

\begin{thmbox}
\textbf{Theorem 7.8 (Characteristic Determined by Unity).} \textit{Let $R$ be a ring with unity $1$.}
\begin{enumerate}[label=\textbf{\arabic*)}]
    \item \textit{If the order of $1$ under addition in $(R, +)$ is finite, say $O(1) = n$, then $\charac(R) = n$.}
    \item \textit{If the order of $1$ under addition is infinite, then $\charac(R) = 0$.}
\end{enumerate}
\end{thmbox}

\begin{proof}
\textbf{Part 1:} Let $O(1) = n$ in $(R, +)$. Then $n$ is the least positive integer such that $n \cdot 1 = 0$. For any arbitrary element $a \in R$:
\[
n \cdot a = \underbrace{a + a + \dots + a}_{n \text{ times}} = \underbrace{(1 + 1 + \dots + 1)}_{n \text{ times}} a = (n \cdot 1)a = 0 \cdot a = 0.
\]
Thus, $n \cdot a = 0$ for all $a \in R$. Furthermore, since $n$ is the least positive integer satisfying $n \cdot 1 = 0$, it is the least positive integer annihilating all elements of $R$. Therefore, $\charac(R) = n$.

\medskip
\noindent \textbf{Part 2:} If $O(1)$ is infinite, then there exists no positive integer $n$ such that $n \cdot 1 = 0$. Consequently, no positive integer can annihilate all elements of $R$, so $\charac(R) = 0$.
\end{proof}

\begin{thmbox}
\textbf{Theorem 7.9 (Characteristic of an Integral Domain).} \textit{The characteristic of an integral domain (and hence of any field) is either $0$ or a prime number.}
\end{thmbox}

\begin{proof}
Let $D$ be an integral domain. If $\charac(D) = 0$, the theorem holds. 

Now suppose $\charac(D) = n > 0$. We must show that $n$ is a prime number. Suppose for contradiction that $n$ is composite. Then there exist integers $n_1, n_2$ such that:
\[
n = n_1 n_2, \quad \text{where } 1 < n_1, n_2 < n.
\]
Since $\charac(D) = n$, by Theorem 7.8 we have $n \cdot 1 = 0$. Therefore:
\[
(n_1 n_2) \cdot 1 = 0 \implies (n_1 \cdot 1)(n_2 \cdot 1) = 0.
\]
Because $D$ is an integral domain, it has no zero divisors. Thus, the above product can only be zero if:
\[
n_1 \cdot 1 = 0 \quad \text{or} \quad n_2 \cdot 1 = 0.
\]
However, this contradicts the definition of $n$ as the \textit{least} positive integer such that $n \cdot 1 = 0$, since both $n_1 < n$ and $n_2 < n$. 

Hence, $n$ cannot be composite. Therefore, $n$ must be a prime number.
\end{proof}

\end{document}
